% !TEX root = ../main_without_quantization.tex

\chapter{Foundations}
% !TEX root = ../../main_without_quantization.tex
\section{Introduction}

The purpose of pruning is to keep the behavior of a trained neural network while modifying its size by removing some of its parameters or biggest parts. In order to study what is eligible to be removed and how the network will be effected by removing these parts it is first needed to explain what are the objects that pruning can work on: Neurons, weights, activations, layers, feed-forward neural network. The vocabulary is not given here as general knowledge.

We begin this chapter by giving a description of a single artificial neuron. This is necessary for pruning because weights and biases are the smallest parameters which are available for training and activation determines where the information should propagate to after weights and biases have been updated. It naturally leads us to multi-layer perceptrons, loss functions, gradient based methods and backpropagation because we measure the performance of a pruning technique by the amount to which the output behavior, or the loss, changes when removing some part of a network\cite{Goodfellow2016,Rumelhart1986}.


\newpage
\section{Parameterization of Neural Networks}
An artificial neuron is a computing unit that takes an input vector, it applies a linear transformation and finally inputs this result to a non-linear activation function. This form of an artificial neuron enables the artificial neural network to learn non-linear mappings. The linear part is also called an affine transformation, which computes a weighted sum of the input vector and then adds a bias. The non-linear activation function is important to enable the artificial neural network to learn complex patterns, since it is responsible for approximate the non-linearities in the data. Examples for common activation functions are ReLU, sigmoid and tanh.

\subsection{The Neuron as a Linear Model}

The artificial neuron, formally introduced by \cite{McCulloch1943}, serves as the fundamental building block of neural networks. Mathematically, it performs an \textbf{affine transformation} on its inputs followed by a non-linear activation $f$ which is represented in Figure~\ref{fig:neuron_affine_unit} and expressed in Eq.\eqref{eq:neuron_affine_transformation}.

\begin{equation}
\label{eq:neuron_affine_transformation}
z = \vect{w}^T \vect{x} + b = \sum_{i=1}^{n} w_i x_i + b
\end{equation}
\noindent Given an input $\vect{x} \in \reals^n$, weights $\vect{w} \in \reals^n$ scale the importance of each feature, while the bias $b \in \reals$ acts as an offset. Figure~\ref{fig:neuron_affine_unit} shows this affine unit before the activation is applied: the weighted inputs are summed with the bias to produce the pre-activation response $z$.

\begin{figure}[htbp]
    \centering
    \includegraphics[width=0.82\textwidth]{neuron_perceptron_with_activation.pdf}
    \caption{A single artificial neuron with input vector $\vect{x}$, weight vector $\vect{w}$, bias $b$, and pre-activation output $z$. The activation function $f$ is applied to $z$ to produce the final output $y$.}
    \label{fig:neuron_affine_unit}
\end{figure}

The term $\vect{w}^T \vect{x}$ measures how strongly the input matches the learned weights. This affine computation is the basic operation repeated across layers in larger neural networks.

% !TEX root = ../../main_without_quantization.tex
\subsection{Activation Functions}

The neural network's ability to model complex patterns depends on the non-linear activation functions applied after each affine transformation. If we were to stack multiple layers without any non-linearity, the entire network would collapse into a single affine transformation.

Therefore, depth alone is not enough: neural networks require non-linear \textbf{activation functions} to model complex patterns \cite{Goodfellow2016}.

As shown in Figure~\ref{fig:neuron_affine_unit}, the activation function $f$ takes the pre-activation output $z$ and produces the final output $y=f(z)$. The choice of activation function can significantly affect the network's performance, convergence, and ability to learn from data.
\subsubsection{Classical Saturating Activations}

Sigmoid (Eq.\eqref{eq:sigmoid}) and hyperbolic tangent (Eq.\eqref{eq:tanh}) functions were often used in early networks. While these functions take real-valued inputs and output bounded values, which could be useful in the case of probabilities or normalized hidden states, they have very small derivatives when the input is very large or very small. This saturation causes learning to be very slow in deep networks \cite{Glorot2010,Goodfellow2016}.

\begin{equation}
    \label{eq:sigmoid}
    \sigma(x)=\frac{1}{1+e^{-x}}, \qquad \sigma:\mathbb{R}\rightarrow(0,1).
\end{equation}

\begin{equation}
    \label{eq:tanh}
    \tanh(x)=\frac{e^x-e^{-x}}{e^x+e^{-x}}, \qquad \tanh:\mathbb{R}\rightarrow(-1,1).
\end{equation}

\subsubsection{Rectified Activations}

Rectified Linear Units (ReLU) became common in deep learning because they are simple, computationally cheap, and reduce saturation for positive inputs \cite{Nair2010}. ReLU is defined as showing in Eq.\eqref{eq:relu}:
\begin{equation}
    \label{eq:relu}
    \operatorname{ReLU}(x)=\max(0,x), \qquad \operatorname{ReLU}:\mathbb{R}\rightarrow[0,\infty).
\end{equation}

Leaky ReLU keeps a small non-zero slope for negative inputs, reducing the risk that a unit becomes permanently inactive \cite{Maas2013}, as shown in Eq.\eqref{eq:leaky_relu}:
\begin{equation}
    \label{eq:leaky_relu}
    \operatorname{LeakyReLU}(x)=
    \begin{cases}
        x, & x \ge 0,\\
        \alpha x, & x<0,
    \end{cases}
    \qquad \alpha \in (0,1), \quad \operatorname{LeakyReLU}:\mathbb{R}\rightarrow\mathbb{R}.
\end{equation}

The Exponential Linear Unit (ELU) uses a smooth negative branch and a linear positive branch \cite{Clevert2016}, which can lead to faster convergence and improved performance in some cases, as shown in Eq.\eqref{eq:elu}:
\begin{equation}
    \label{eq:elu}
    \operatorname{ELU}(x)=
    \begin{cases}
        x, & x \ge 0,\\
        \alpha(e^x-1), & x<0,
    \end{cases}
    \qquad \alpha>0, \quad \operatorname{ELU}:\mathbb{R}\rightarrow(-\alpha,\infty).
\end{equation}

Modern Transformer architectures often use the Gaussian Error Linear Unit (GELU), a smooth activation that weights inputs by their value under the standard Gaussian cumulative distribution \cite{Hendrycks2016}, as shown in Eq.\eqref{eq:gelu}:
\begin{equation}
    \label{eq:gelu}
    \operatorname{GELU}(x)=x\Phi(x)\approx \frac{x}{2}\left(1+\tanh\left(\sqrt{\frac{2}{\pi}}(x+0.044715x^3)\right)\right).
\end{equation}

Figure~\ref{fig:activation_comparison} compares the shapes of these activation functions on the same input range. The plot makes the difference between saturating functions and rectified functions visible.

\begin{figure}[H]
    \centering
    \begin{tikzpicture}
        \begin{axis}[
            width=0.82\textwidth,
            height=6.2cm,
            xmin=-4, xmax=4,
            ymin=-1.4, ymax=4,
            axis lines=middle,
            xlabel={$x$},
            ylabel={$f(x)$},
            samples=200,
            grid=major,
            legend style={at={(0.02,0.98)},anchor=north west,font=\small},
        ]
            \addplot[deepmutedblue, very thick] {1/(1+exp(-x))};
            \addlegendentry{Sigmoid}
            \addplot[dustyteal, very thick] {tanh(x)};
            \addlegendentry{Tanh}
            \addplot[orange!85!black, very thick] {max(0,x)};
            \addlegendentry{ReLU}
            \addplot[slateblue, thick, dashed] {x < 0 ? 0.1*x : x};
            \addlegendentry{Leaky ReLU}
            \addplot[actreddark, thick, dashdotted] {x < 0 ? exp(x)-1 : x};
            \addlegendentry{ELU}
            \addplot[black, thick, dotted] {0.5*x*(1+tanh(sqrt(2/pi)*(x+0.044715*x^3)))};
            \addlegendentry{GELU}
        \end{axis}
    \end{tikzpicture}
    \caption{Comparison of common activation functions}
    \label{fig:activation_comparison}
\end{figure}

% !TEX root = ../../main_without_quantization.tex

\subsection{The Multi-Layer Perceptron (MLP)}

A single activated neuron can transform an input into a simple response, but practical models require many such units to work together. A Multi-Layer Perceptron extends the neuron model by organizing neurons into successive layers, so that each layer transforms the representation produced by the one before it.

Going from a single neuron to an MLP increases expressiveness by composing affine transformations with non-linear activations across layers~\cite{Rumelhart1986}. MLPs are \emph{feedforward} and typically \emph{fully connected}, meaning each neuron connects to all neurons in the next layer.

Eq.~\eqref{eq:mlp_layer_computation} formalizes the computation performed by each layer of an MLP, to transform the raw input $\vect{x}$ into the final output (prediction)$\vect{y}$.
\begin{equation}
\label{eq:mlp_layer_computation}
\begin{aligned}
	\vect{h}^{(0)} &= \vect{x},\\
	\vect{h}^{(\ell)} &= f(\matr{W}^{(\ell)}\vect{h}^{(\ell-1)} + \vect{b}^{(\ell)}) \quad \text{for } \ell = 1, 2, \ldots, L-1,\\
	\vect{y} &= f_{\text{out}}(\matr{W}^{(L)}\vect{h}^{(L-1)} + \vect{b}^{(L)}),
\end{aligned}
\end{equation}

	Figure~\ref{fig:mlp_architecture} illustrates a typical MLP with one input layer, two hidden layers, and one output layer.
\begin{figure}[ht]
    \centering
    \resizebox{0.82\textwidth}{!}{%
    \begin{tikzpicture}[
        scale=1.0,
        node distance=2.5cm,
        neuron/.style={circle, draw=black!70, fill=black!8, line width=1.4pt, minimum size=0.7cm},
        input/.style={circle, draw=black!85, fill=black!15, line width=1.6pt, minimum size=0.7cm},
        output/.style={circle, draw=black!70, fill=black!5, line width=1.4pt, minimum size=0.7cm},
        arrow/.style={-Stealth, color=black!30, thin},
        highlight/.style={-Stealth, color=black!65, very thick},
        annotation/.style={font=\scriptsize, color=black!65}
        ]

        % Input layer
        \foreach \i in {1,2,3,4} {
            \node[input] (x\i) at (0, -\i*1.2) {};   % was 1.5
        }
        \node[above=0.3cm of x1, color=black!75, font=\bfseries] {Input Layer};
        \node[right=0.2cm of x1, color=black!60, font=\small] {$x_1$};
        \node[right=0.2cm of x2, color=black!60, font=\small] {$x_2$};
        \node[right=0.2cm of x3, color=black!60, font=\small] {$x_3$};
        \node[right=0.2cm of x4, color=black!60, font=\small] {$x_4$};
        \node[below=0.3cm of x4, color=black!60, font=\small] {$\vect{h}^{(0)} \in \reals^{4}$};

        % Hidden layer 1
        \foreach \i in {1,2,3,4,5} {
            \node[neuron] (h1\i) at (4, -\i*1.1) {};   % was 1.3
        }
        \node[above=0.3cm of h11, color=black!75, font=\bfseries] {Hidden Layer 1};
        \node[below=0.3cm of h15, color=black!60, font=\small] {$\vect{h}^{(1)} \in \reals^{5}$};

        % Hidden layer 2
        \foreach \i in {1,2,3,4} {
            \node[neuron] (h2\i) at (8, -\i*1.2) {};   % was 1.5
        }
        \node[above=0.3cm of h21, color=black!75, font=\bfseries] {Hidden Layer 2};
        \node[below=0.3cm of h24, color=black!60, font=\small] {$\vect{h}^{(2)} \in \reals^{4}$};

        % Output layer
        \foreach \i in {1,2} {
            \node[output] (y\i) at (12, {-1.8-(\i-1)*1.8}) {};   % tightened gap
        }
        \node[above=0.3cm of y1, color=black!75, font=\bfseries] {Output Layer};
        \node[left=0.2cm of y1, color=black!60, font=\small] {$y_1$};
        \node[left=0.2cm of y2, color=black!60, font=\small] {$y_2$};
        \node[below=0.3cm of y2, color=black!60, font=\small] {$\vect{y} \in \reals^{2}$};

        % Connections: Input to Hidden 1
        \foreach \i in {1,2,3,4} {
            \foreach \j in {1,2,3,4,5} {
                \draw[arrow] (x\i) -- (h1\j);
            }
        }
        % Connections: Hidden 1 to Hidden 2
        \foreach \i in {1,2,3,4,5} {
            \foreach \j in {1,2,3,4} {
                \draw[arrow] (h1\i) -- (h2\j);
            }
        }
        % Connections: Hidden 2 to Output
        \foreach \i in {1,2,3,4} {
            \foreach \j in {1,2} {
                \draw[arrow] (h2\i) -- (y\j);
            }
        }

        % Highlight one example path
        \draw[highlight] (x2) -- (h13);
        \draw[highlight] (h13) -- (h22);
        \draw[highlight] (h22) -- (y1);

        % Weight matrix labels
        \node[color=black!65, font=\small\bfseries, align=center] at (2, -7.2) {$\matr{W}^{(1)}$\\[0.1cm] \scriptsize $(5 \times 4)$};
        \node[color=black!65, font=\small\bfseries, align=center] at (6, -7.2) {$\matr{W}^{(2)}$\\[0.1cm] \scriptsize $(4 \times 5)$};
        \node[color=black!65, font=\small\bfseries, align=center] at (10, -7.2) {$\matr{W}^{(3)}$\\[0.1cm] \scriptsize $(2 \times 4)$};

        % Bias vectors
        \node[annotation, align=center] at (2, -8.0) {$\vect{b}^{(1)} \in \reals^5$};
        \node[annotation, align=center] at (6, -8.0) {$\vect{b}^{(2)} \in \reals^4$};
        \node[annotation, align=center] at (10, -8.0) {$\vect{b}^{(3)} \in \reals^2$};

        % Activation annotations
        \node[annotation, align=center] at (2, 0.6) {$\vect{h}^{(1)} = \sigma(\matr{W}^{(1)}\vect{h}^{(0)} + \vect{b}^{(1)})$};
        \node[annotation, align=center] at (6, 0.6) {$\vect{h}^{(2)} = \sigma(\matr{W}^{(2)}\vect{h}^{(1)} + \vect{b}^{(2)})$};
        \node[annotation, align=center] at (10, 0.6) {$\vect{y} = \sigma(\matr{W}^{(3)}\vect{h}^{(2)} + \vect{b}^{(3)})$};

        % Legend
        \node[draw=black!30, fill=white, rounded corners, align=left, font=\scriptsize] at (15, -3) {
            \textbf{Legend:}\\[0.2cm]
            \tikz\draw[arrow] (0,0) -- (0.4,0); Weighted connection\\[0.15cm]
            \tikz\draw[highlight] (0,0) -- (0.4,0); Example path\\[0.15cm]
            \tikz\node[input, minimum size=0.4cm] at (0.2,0) {}; Input neuron\\[0.15cm]
            \tikz\node[neuron, minimum size=0.4cm] at (0.2,0) {}; Hidden neuron\\[0.15cm]
            \tikz\node[output, minimum size=0.4cm] at (0.2,0) {}; Output neuron
        };

        % Parameter count box
        \node[draw=black!30, fill=black!5, rounded corners, align=center, font=\scriptsize] at (6, -9.0) {
            \textbf{Total Parameters:} $(5\times4) + 5 + (4\times5) + 4 + (2\times4) + 2 = 59$
        };

        % Top annotation
        \node[annotation, font=\small\bfseries] at (6, 1.4) {Layer-wise composition of affine maps and activations};

    \end{tikzpicture}
    }
    \caption{Fully connected MLP with 4 input features, two hidden layers (5 and 4 neurons), and 2 output neurons.}
    \label{fig:mlp_architecture}
\end{figure}
\newpage
%The diagram shows the layer-wise transformations that define an MLP. The input layer stores the feature vector $\vect{h}^{(0)}=\vect{x}$. The first hidden layer computes $\vect{h}^{(1)}$ from all input features using $\matr{W}^{(1)}$ and $\vect{b}^{(1)}$, then applies a non-linearity. The second hidden layer repeats the same operation on $\vect{h}^{(1)}$, producing $\vect{h}^{(2)}$. Finally, the output layer maps the last hidden representation to the prediction vector $\vect{y}$. The highlighted orange path illustrates one route by which information from a single input feature can influence an output neuron, and the pale connections show that every layer is fully connected to the next.

%\subsection{Hidden Representations}
%
%The hidden layers of an MLP are its internal states. The $i$-th hidden layer effectively performs the $i$-th transformation of the data, mapping it into the $i$-th coordinate space discovered through training. The earlier hidden layers perform combinations of input features, while the later layers perform combinations of the earlier layer outputs, creating higher-level representations related to the task in question. The output of the last hidden layer is typically transformed into the appropriate output space. For classification, this means class scores or probabilities. For regression, it means one or more continuous values.
%
%The multi-layer representation is helpful, but it suffers from an obvious limitation. The MLP accepts its input as a single vector; therefore, there is no innate reason in the model for elements close to each other to be related, or for the elements of the vector to have an order. If such structure is required, it must be supplied by the architecture, which we will tackle later in Section~\ref{sec:specialized_architectures}.
%
%For now, we remain with the MLP and ask a more basic question.
%Once the architecture has fixed how information can flow from inputs to outputs, how are the actual parameters of this flow chosen? This is the role of learning: the weights and biases must be adjusted from data so that the network output becomes closer to the desired target. The next section \ref{sec:Optimization and Gradients}
%formulates this process as an optimization problem.

% !TEX root = ../../main_without_quantization.tex
\section{Optimization and Gradients}\label{sec:Optimization and Gradients}

We outlined how an MLP generates a prediction: by a composition of transformations applied at each layer. With forward computation, learning can be thought of as figuring out what the weights and biases should be given the data. Network construction then becomes an optimization problem: select the parameters that make the output as close as possible to the target on a training set.

Hence, a quantitative measure must be provided that can actually be minimized; this provides the rationale for empirical risk minimization.

\subsection{Empirical Risk Minimization}

Given a dataset $\mathcal{D} = \{(\vect{x}^{(i)}, y^{(i)})\}_{i=1}^{m}$, learning consists of finding parameters $\params$ that minimize the discrepancy between the model predictions $f_{\params}(\vect{x})$ and the target values. This is commonly formulated as empirical risk minimization \cite{Vapnik1998}:
\begin{equation}
    \params^* = \arg\min_{\params} \hat{R}(\params) = \arg\min_{\params} \frac{1}{m} \sum_{i=1}^{m} \loss(f_{\params}(\vect{x}^{(i)}), y^{(i)}).
    \label{eq:erm}
\end{equation}
where $m$ is the number of training samples; $\vect{x}^{(i)}$ and $y^{(i)}$ are the input and target of the $i$-th sample; $\params$ are the model parameters; $\hat{R}(\params)$ is the empirical risk, i.e.\ the average loss over the training set; and $\loss(\cdot,\cdot)$ is the per-sample loss function. For a $K$-class classification problem the prediction is a probability vector $\hat{\vect{y}} \in \Delta^{K}$ (the probability simplex over $K$ classes) with entries $\hat{y}_k$, the target $\vect{y}$ has entries $y_k$, and $c$ denotes the index of the true class; the Kullback--Leibler divergence ~\cite{cover2006}
$D_{\text{KL}}(\vect{y} \,\|\, \hat{\vect{y}})$ measures the excess cost of using $\hat{\vect{y}}$ in place of $\vect{y}$.

This empirical-risk objective in Eq.~\eqref{eq:erm} gives the general learning template: choose parameters by minimizing the average loss over the training set. It does not specify the loss itself. That choice is determined by the problem and the structure of the output.

\subsubsection{Loss Functions}

Now we have an optimization problem. But what are we minimizing? A loss function accounts for the discrepancy between predictions and targets and makes explicit assumptions about the problem and the output space. Table~\ref{tab:loss_functions} summarizes the most common loss functions used in neural network training, along with their formulas, target types, and properties. The choice of loss function depends on the task at hand (classification vs regression) and the nature of the targets (hard one-hot labels, soft probability distributions, or continuous real values.).
%\subsubsection{Cross-Entropy Loss}
%
%For classification tasks, neural networks typically produce a probability distribution over $K$ classes via a softmax output. Given a predicted distribution $\hat{\vect{y}} \in \Delta^{K}$ and a target distribution $\vect{y}$ (often a one-hot vector), the cross-entropy loss is defined as:
%\begin{equation}
%    \loss_{\text{CE}}(\hat{\vect{y}}, \vect{y}) = -\sum_{k=1}^{K} y_k \log \hat{y}_k.
%    \label{eq:ce}
%\end{equation}
%The cross-entropy expression in Equation~\eqref{eq:ce} quantifies the number of bits required on average to encode symbols from $\vect{y}$ using a code optimal for $\hat{\vect{y}}$. For a hard one-hot label, this becomes the negative log-likelihood of the true class $\loss_{\text{CE}} = -\log \hat{y}_c$, where $c$ is the index of the true class. This has the effect of pushing most of the probability mass onto the correct class, and preventing it from going anywhere else. For a hard label, this works out fine, but the moment the target becomes a distribution instead of a single class, the problem becomes more explicitly about how to quantify distance between distributions, which motivates the following point.
%
%\subsubsection{Kullback--Leibler Divergence}
%
%Cross-entropy implicitly compares two distributions, which raises a natural question: what is the right way to measure how far apart two probability distributions are? The Kullback--Leibler (KL) divergence \cite{Kullback1951} provides a principled answer rooted in information theory. Rather than measuring total encoding cost, it isolates the \emph{excess} cost incurred by using $\hat{\vect{y}}$ in place of the true distribution $\vect{y}$:
%\begin{equation}
%    D_{\text{KL}}(\vect{y} \,\|\, \hat{\vect{y}}) = \sum_{k=1}^{K} y_k \log \frac{y_k}{\hat{y}_k}.
%    \label{eq:kl}
%\end{equation}
%Expanding the KL divergence in Equation~\eqref{eq:kl} reveals the exact relationship with cross-entropy:
%\begin{equation}
%    D_{\text{KL}}(\vect{y} \,\|\, \hat{\vect{y}})
%    = \underbrace{-\sum_{k} y_k \log \hat{y}_k}_{\loss_{\text{CE}}(\hat{\vect{y}},\,\vect{y})}
%    \;-\; \underbrace{\left(-\sum_{k} y_k \log y_k\right)}_{H(\vect{y})},
%    \label{eq:kl_ce_decomp}
%\end{equation}
%where $H(\vect{y})$ is the entropy of the target distribution. The decomposition in Equation~\eqref{eq:kl_ce_decomp} shows why cross-entropy and KL divergence lead to the same optimization problem when the target distribution is fixed: $H(\vect{y})$ does not depend on the model parameters. If the target is a one-hot label then $H(\vect{y}) = 0$ and the two functions are equal; if the target is a soft distribution, then $H(\vect{y}) > 0$, and KL divergence is the clearly superior objective. The above loss functions are for when the output is a discrete probability distribution. For real-valued regression problems, an alternative family of loss functions are used.
%
%\subsubsection{Regression Losses}
%
%For regression, the output is a continuous value rather than a class index, and the appropriate loss changes accordingly. The mean squared error (MSE) is the standard choice:
%\begin{equation}
%    \loss_{\text{MSE}}(\hat{y}, y) = \frac{1}{m}\sum_{i=1}^{m}(\hat{y}^{(i)} - y^{(i)})^2.
%    \label{eq:mse}
%\end{equation}
%The MSE objective in Equation~\eqref{eq:mse} penalizes large regression errors quadratically. Under a Gaussian noise assumption on the targets, minimizing it is equivalent to maximum likelihood estimation of the model parameters. When the output distribution is heavier-tailed or contains outliers, the mean absolute error (L1 loss) $\loss_{\text{MAE}} = \frac{1}{m}\sum_i |\hat{y}^{(i)} - y^{(i)}|$ is more robust, as it does not square large deviations.
%
%\noindent\rule{0.4\linewidth}{0.55pt}
\begin{table}[ht]
\centering
\caption{Summary of common neural network loss functions.}
\label{tab:loss_functions}
\renewcommand{\arraystretch}{1.45}
\begin{tabularx}{\textwidth}{@{} l l X l X @{}}
\toprule
\textbf{Loss} & \textbf{Task} & \textbf{Formula} & \textbf{Target type} & \textbf{properties} \\
\midrule

Cross-Entropy
  & Classification
  & $\loss_{\text{CE}}(\hat{\vect{y}}, \vect{y}) = -\displaystyle\sum_{k=1}^{K} y_k \log \hat{y}_k$
  & Soft or one-hot
  & Measures total encoding cost of $\vect{y}$ under code optimal for $\hat{\vect{y}}$;
    reduces to $-\log\hat{y}_c$ for hard labels. \\

KL Divergence
  & Classification
  & $D_{\text{KL}}(\vect{y}\,\|\,\hat{\vect{y}}) = \displaystyle\sum_{k=1}^{K} y_k \log \dfrac{y_k}{\hat{y}_k}$
  & Soft distribution
  & Measures \emph{excess} encoding cost;
    $D_{\text{KL}} = \loss_{\text{CE}} - H(\vect{y})$.
    Preferred over CE when targets are soft. \\

MSE (L2)
  & Regression
  & $\loss_{\text{MSE}}(\hat{y}, y) = \dfrac{1}{m}\displaystyle\sum_{i=1}^{m}(\hat{y}^{(i)} - y^{(i)})^2$
  & Continuous
  & Penalises errors quadratically; equivalent to maximum likelihood estimation (MLE) under Gaussian noise.
    Sensitive to outliers. \\

MAE (L1)
  & Regression
  & $\loss_{\text{MAE}}(\hat{y}, y) = \dfrac{1}{m}\displaystyle\sum_{i=1}^{m}|\hat{y}^{(i)} - y^{(i)}|$
  & Continuous
  & Penalises errors linearly; robust to heavy-tailed distributions and outliers. \\

\bottomrule
\end{tabularx}

\smallskip
\begin{tabularx}{\textwidth}{@{} X @{}}
\small
\textbf{Notes.}
$K$: number of classes;
$\hat{\vect{y}}\in\Delta^{K}$: predicted probability vector (softmax output);
$\vect{y}$: target distribution (one-hot or soft);
$c$: index of the true class;
$H(\vect{y}) = -\sum_k y_k\log y_k$: entropy of the target;
$m$: number of samples;
$\hat{y}^{(i)}, y^{(i)}$: predicted and true scalar outputs for sample $i$.
When $\vect{y}$ is one-hot, $H(\vect{y})=0$ and $D_{\text{KL}}=\loss_{\text{CE}}$.
\\
\end{tabularx}
\end{table}

% !TEX root = ../../main_without_quantization.tex
\subsection{Algorithms for Gradient-Based Optimization}
\label{sec:gradient_based_optimization}
With the loss function established, the remaining question is how the resulting error signal propagates back through the network to update the parameters.

Training neural networks is usually a problem of finding a parameter vector $\theta \in \mathbb{R}^{d}$ that minimizes the scalar objective $J(\theta)$ where $J$ is the empirical risk defined by the loss function and the training data. Formally, we want to solve \eqref{eq:optimization_problem}:
\begin{equation}
    \theta^* = \arg\min_{\theta} J(\theta).
    \label{eq:optimization_problem}
\end{equation}

The parameter space in most contemporary neural networks is high-dimensional and the objective is largely non-convex so a solution in closed form for optimum parameters rarely exists. Training is hence iterative: the parameters are updated according to directions suggested by the gradient of the loss \cite{Bottou2018}.

The gradient $\nabla J(\theta)$ indicates the direction of steepest local increase of the objective, and optimizers will use it in reverse to adjust the parameters to  yield a lower loss. This concept forms the basis of the gradient descent algorithm as well as of all the adaptive optimization schemes currently used in training neural networks.

\subsubsection{Gradient Descent}

Gradient descent updates parameters in the negative-gradient direction. For iteration $t$ and learning rate $\eta>0$, the update is formulated in Eq.\eqref{eq:gradient_descent}:
\begin{equation}
    \label{eq:gradient_descent}
    \theta_{t+1} = \theta_t - \eta \nabla J(\theta_t).
\end{equation}

The learning rate controls the step size. If it is too small, optimization may progress slowly; if it is too large, the updates can overshoot useful regions of the objective surface. In large scale learning, the exact gradient over the full dataset is often replaced by a stochastic estimate computed on a mini-batch. This gives stochastic gradient descent, which is cheaper per update and is widely used for neural-network training \cite{Bottou2018}.

\subsubsection{Adam}

Adam is an adaptive gradient-based optimizer that combines momentum with per-parameter learning-rate adaptation \cite{kingma2015}. Instead of using only the current gradient, Adam keeps exponential moving averages of the first and second moments of the gradient, which are updated as shown in Eq.\eqref{eq:adam_moments}:
\begin {equation}
    \label{eq:adam_moments}
    m_t = \beta_1 m_{t-1} + (1-\beta_1)g_t,
    \qquad
    v_t = \beta_2 v_{t-1} + (1-\beta_2)g_t^2,
\end{equation}

%\begin{align}
%    m_t &= \beta_1 m_{t-1} + (1-\beta_1)g_t,\\
%    v_t &= \beta_2 v_{t-1} + (1-\beta_2)g_t^2,
%\end{align}
where $g_t=\nabla J(\theta_t)$ is the gradient estimate at iteration $t$, $m_t$ tracks the average direction of recent gradients, and $v_t$ tracks their squared magnitude.

Because both moving averages are initialized at zero, Adam applies bias correction computed as follows Eq.\eqref{eq:adam_bias_correction}:
\begin{equation}
    \label{eq:adam_bias_correction}
    \hat{m}_t = \frac{m_t}{1-\beta_1^t},
    \qquad
    \hat{v}_t = \frac{v_t}{1-\beta_2^t}.
\end{equation}
The parameter update is then given by \eqref{eq:adam_update}:
\begin{equation}
    \label{eq:adam_update}
    \theta_{t+1}
    =
    \theta_t
    -
    \eta \frac{\hat{m}_t}{\sqrt{\hat{v}_t}+\epsilon}.
\end{equation}

Adam is often utilized due to its robustness in relation to scaling differences among gradients for different parameters and typically requires less fine-tuning of learning rates as opposed to simple gradient descent. Subsequent versions of Adam include AdamW, which removes the weight decay from the adaptive learning component and is applicable to a variety of deep learning applications \cite{Loshchilov2019}.



Both gradient descent and Adam describe how the parameters should be updated, but they both require that the gradients be available. These computations can become expensive in deep, multi-layered networks and thus require an alternative method, termed backpropagation.
% !TEX root = ../../main_without_quantization.tex
\section{Backpropagation}

As established in Section~\ref{sec:gradient_based_optimization}, gradient descent and Adam require gradients of the loss with respect to the parameters. Backpropagation is the mechanism that connects the loss at the output of the network to the parameters inside each layer by applying the chain rule backward through the network \cite{Rumelhart1986}.

During the forward pass, each layer transforms the representation produced by the previous one. If the input to layer $l$ is $a^{(l-1)}$, then the layer first forms
\begin{align}
    z^{(l)} &= W^{(l)}a^{(l-1)} + b^{(l)},\\
    a^{(l)} &= \phi\!\left(z^{(l)}\right),
\end{align}
where $W^{(l)}$ is the weight matrix, $b^{(l)}$ is the bias vector, $z^{(l)}$ is the pre-activation, $a^{(l)}$ is the activated output, and $\phi$ is the activation function. The loss $\mathcal{L}(a^{(L)},y)$ is only observed at the final output, but every earlier layer contributed to that output through the sequence of transformations that followed it.

The central idea is therefore to move the effect of the loss backward through the network. At each layer, backpropagation keeps track of how sensitive the loss is to the pre-activation value of that layer:
\begin{equation}
    \delta^{(l)} = \frac{\partial \mathcal{L}}{\partial z^{(l)}}.
\end{equation}
For the last layer, this quantity is obtained directly from the derivative of the loss and the derivative of the activation:
\begin{equation}
    \delta^{(L)}
    =
    \frac{\partial \mathcal{L}}{\partial a^{(L)}}
    \odot
    \phi'\!\left(z^{(L)}\right),
\end{equation}
where $\odot$ denotes element-wise multiplication. For an earlier layer, the signal is obtained from the layer after it:
\begin{equation}
    \delta^{(l)}
    =
    \left(W^{(l+1)}\right)^{T}\delta^{(l+1)}
    \odot
    \phi'\!\left(z^{(l)}\right).
\end{equation}
This is the chain rule applied repeatedly along the computational path from the loss back to the parameters.

Once this backward signal is known, the gradients of the layer parameters follow from the same local quantities used in the forward pass:
\begin{align}
    \frac{\partial \mathcal{L}}{\partial W^{(l)}} &= \delta^{(l)}\left(a^{(l-1)}\right)^{T},\\
    \frac{\partial \mathcal{L}}{\partial b^{(l)}} &= \delta^{(l)}.
\end{align}
The update of a weight is therefore tied to the activation that entered the layer and to the error signal that returned to it. Backpropagation avoids deriving each parameter derivative separately; it reuses the stored forward activations and propagates one backward signal from layer to layer.

\vspace{0.8em}
\noindent\rule{\textwidth}{0.4pt}
\vspace{0.8em}

Since prediction is done by forward propagation and learning is done by backpropagation, we have established the core mechanism of training a neural network. This mechanism does not depend on a specific architecture: if the network defines a differentiable function from inputs to outputs, its parameters can be trained by this method. What differs between network families is the set of assumptions they make about the input data.

This is the idea behind the following section, which breaks away from the general MLP setup and looks at specific architectures. The rationale for specific architectures is that some inputs cannot be treated as abstract, unordered vectors, as is the case with time series because of temporal order, compressed representations because of latent structure, and language because of word dependencies.

% !TEX root = ../../main_without_quantization.tex
\section{Specialized Architectures}
\label{sec:specialized_architectures}

While Multi-Layer Perceptrons offer a generic learning model over vector inputs, they do not inherently leverage any underlying structure of the data, and thus can only indirectly learn the temporal dependencies of sequences, the latent structure of data compressed in autoencoders, or the relational structure of positions within a context in token-based models. This fundamental limitation gives rise to architectures tailored to the structure of the data.

\subsection{Recurrent Neural Networks (RNNs)}
RNNs are designed for sequential data (e.g., time series, text). Unlike feedforward networks, RNNs utilize cyclic connections to maintain a \textit{hidden state} that captures historical information. The network dynamics are described by:
\[
h_t = \sigma(W_{hh} h_{t-1} + W_{xh} x_t + b_h)
\]
\[
y_t = W_{hy} h_t + b_y
\]
where \( h_t \) is the hidden state, \( x_t \) the input, and \( y_t \) the output at time \( t \).
Figure~\ref{fig:rnn-compact} represents this recurrence by unrolling the same recurrent computation across time steps, so that the dependence of $h_t$ on $h_{t-1}$ becomes visible.
\begin{figure}[htbp]
  \centering
  \begin{tikzpicture}[>=Stealth, node distance=10mm and 18mm, every node/.style={font=\small}]
    \tikzstyle{inp} = [circle, draw=slateblue, fill=mistyblue!20, minimum size=7mm]
    \tikzstyle{hid} = [rectangle, draw=dustyteal, fill=sagegreen!20, minimum size=9mm]
    \tikzstyle{outnode} = [circle, draw=slateblue, fill=mistyblue!20, minimum size=7mm]
    \tikzstyle{lbl} = [font=\scriptsize, color=slateblue]
    
    % Inputs (bottom layer)
    \node[inp] (x1) {$x_1$};
    \node[inp] (x2) [right=of x1] {$x_2$};
    \node[inp] (x3) [right=of x2] {$x_3$};
    \node (xd) [right=of x3] {$\dots$};
    \node[inp] (xT) [right=of xd] {$x_T$};
    
    % Hidden (middle layer)
    \node[hid] (h1) [above=of x1] {$h_1$};
    \node[hid] (h2) [above=of x2] {$h_2$};
    \node[hid] (h3) [above=of x3] {$h_3$};
    \node[hid] (hT) [above=of xT] {$h_T$};
    \node (hd) [right=of h3] {$\dots$};
    
    % Outputs (top layer)
    \node[outnode] (y1) [above=of h1] {$y_1$};
    \node[outnode] (y2) [above=of h2] {$y_2$};
    \node[outnode] (y3) [above=of h3] {$y_3$};
    \node[outnode] (yT) [above=of hT] {$y_T$};
    
    % Input to hidden
    \draw[->] (x1) -- node[midway, left, lbl] {$W_{xh}$} (h1);
    \draw[->] (x2) -- node[midway, left, lbl] {$W_{xh}$} (h2);
    \draw[->] (x3) -- node[midway, left, lbl] {$W_{xh}$} (h3);
    \draw[->] (xT) -- node[midway, left, lbl] {$W_{xh}$} (hT);
    
    % Hidden recurrence 
    \draw[->] (h1) to[bend left=8] node[midway, above, lbl] {$W_{hh}$} (h2);
    \draw[->] (h2) to[bend left=8] node[midway, above, lbl] {$W_{hh}$} (h3);
    \draw[->] (h3) -- (hd);
    \draw[->] (hd) to[bend left=8] node[midway, above, lbl] {$W_{hh}$} (hT);
    
    % Hidden to Output
    \draw[->] (h1) -- node[midway, left, lbl] {$W_{hy}$} (y1);
    \draw[->] (h2) -- node[midway, left, lbl] {$W_{hy}$} (y2);
    \draw[->] (h3) -- node[midway, left, lbl] {$W_{hy}$} (y3);
    \draw[->] (hT) -- node[midway, left, lbl] {$W_{hy}$} (yT);
    
    \node[below left=2mm and 6mm of h1, font=\scriptsize] (h0label) {$h_0$};
    \draw[->] (h0label.north) to[bend left=20] (h1.west);
  \end{tikzpicture}
  \caption{Recurrent Neural Network (RNN) unrolled in time. The input $x_t$ at each time step updates the hidden state $h_t$ through the recurrent weights $W_{hh}$, while $W_{xh}$ maps inputs to hidden states and $W_{hy}$ maps hidden states to outputs $y_t$.}
  \label{fig:rnn-compact}
\end{figure}

RNNs focus on preserving information across time.
Another important architectural idea is to preserve only the most essential information by compressing the input into a lower-dimensional representation.

\subsection{Autoencoders (AEs)}

Autoencoders are neural networks designed for unsupervised learning of efficient data codings. They act as an identity function that learns to compress and then reconstruct data. Figure~\ref{fig:autoencoder} shows this structure as an encoder that maps the input to a latent space and a decoder that reconstructs the output from that compressed representation.

\begin{figure}[htbp]
    \centering
    \begin{tikzpicture}[>=Stealth, node distance=15mm, every node/.style={font=\small}]
        
        \tikzstyle{neuron} = [circle, draw=slateblue, fill=mistyblue!20, minimum size=7mm, inner sep=0pt]
        \tikzstyle{bottleneck} = [circle, draw=dustyteal, fill=sagegreen!20, minimum size=7mm, inner sep=0pt]
        \tikzstyle{lbl} = [font=\scriptsize, color=slateblue]

        % --- Nodes ---
        
        % Input Layer (x)
        \node[neuron] (i1) at (0, 2.0) {};
        \node[neuron] (i2) at (0, 1.0) {};
        \node[neuron] (i3) at (0, 0.0) {};
        \node[neuron] (i4) at (0, -1.0) {};
        \node[neuron] (i5) at (0, -2.0) {};
        
        % Encoder Hidden
        \node[neuron] (e1) at (2, 1.2) {};
        \node[neuron] (e2) at (2, 0) {};
        \node[neuron] (e3) at (2, -1.2) {};

        % Bottleneck / Latent Space (z)
        \node[bottleneck] (z1) at (4, 0.5) {};
        \node[bottleneck] (z2) at (4, -0.5) {};

        % Decoder Hidden
        \node[neuron] (d1) at (6, 1.2) {};
        \node[neuron] (d2) at (6, 0) {};
        \node[neuron] (d3) at (6, -1.2) {};

        % Output Layer (x_hat)
        \node[neuron] (o1) at (8, 2.0) {};
        \node[neuron] (o2) at (8, 1.0) {};
        \node[neuron] (o3) at (8, 0.0) {};
        \node[neuron] (o4) at (8, -1.0) {};
        \node[neuron] (o5) at (8, -2.0) {};

        % --- Connections ---
        
        % Input -> Encoder
        \foreach \i in {1,...,5}
            \foreach \j in {1,...,3}
                \draw[->] (i\i) -- (e\j);

        % Encoder -> Bottleneck
        \foreach \i in {1,...,3}
            \foreach \j in {1,...,2}
                \draw[->] (e\i) -- (z\j);
                
        % Bottleneck -> Decoder
        \foreach \i in {1,...,2}
            \foreach \j in {1,...,3}
                \draw[->] (z\i) -- (d\j);

        % Decoder -> Output
        \foreach \i in {1,...,3}
            \foreach \j in {1,...,5}
                \draw[->] (d\i) -- (o\j);

        % --- Labels ---
        
        \node[below=0.2cm of i5] {Input $x$};
        \node[below=0.2cm of o5] {Output $\hat{x}$};
        \node[below=0.2cm of z2, align=center] {Latent\\Space $z$};
        
        % Braces for Encoder/Decoder
        \draw[decorate, decoration={brace, amplitude=10pt, raise=5pt}, thick, slateblue] (i1.north) -- (z1.north) node[midway, above=20pt] {\textbf{Encoder}};
        \draw[decorate, decoration={brace, amplitude=10pt, raise=5pt}, thick, slateblue] (z1.north) -- (o1.north) node[midway, above=20pt] {\textbf{Decoder}};

    \end{tikzpicture}
    \caption[Autoencoder architecture]{Autoencoder architecture. The encoder compresses the input $x$ into a lower-dimensional latent space $z$, also called the bottleneck, and the decoder reconstructs the output $\hat{x}$ from that latent representation.}
    \label{fig:autoencoder}
\end{figure}

\begin{itemize}
    \item \textbf{Encoder:} Maps the high-dimensional input $x$ to a lower-dimensional latent representation $z$.
    \item \textbf{Decoder:} Reconstructs the output $\hat{x}$ from the latent representation $z$.
\end{itemize}

The network is trained to minimize the reconstruction error (e.g., Mean Squared Error) between the input and the output:
\[
L(x, \hat{x}) = || x - \hat{x} ||^2
\]
By forcing data through a bottleneck, the network learns to capture the most salient features of the data distribution.

Autoencoders show how architectural structure can impose a useful compression constraint.
For the language and sequence models considered later in this report, the most important specialized architecture is the Transformer.




\subsection{Transformer Networks}

Transformer networks, introduced by \cite{vaswani2017attention}, changed sequence modeling by replacing recurrent processing with attention-based computation. Instead of processing tokens one after another, a Transformer can compare tokens across the whole sequence in parallel. This makes it highly effective for language modelling, translation, summarization, and other sequence tasks. The original architecture was an encoder-decoder model for machine translation, while later systems often use encoder-only variants such as BERT \cite{devlin2019bert} or decoder-only variants such as LLaMA \cite{touvron2023llama2}.

\subsection{Core Architectural Components: Building Blocks of the Transformer}

Before separating the model into encoder and decoder stacks, it is useful to define the components repeated inside those stacks. Transformer layers are built from token embeddings, positional information, attention projections, feed-forward channels, residual connections, and normalization layers. These components are also the main targets for compression, because they determine both parameter count and runtime cost.

\subsubsection{Input Representation: Embeddings and Positional Encoding}

Every token is mapped to a continuous embedding vector. Since self-attention does not inherently encode token order, positional information is added to the token representations. The original Transformer uses sinusoidal positional encodings \cite{vaswani2017attention}:
\begin{equation}
    PE_{(\text{pos}, 2i)} = \sin\left(\frac{\text{pos}}{10000^{2i/d_{\text{model}}}}\right), \quad
    PE_{(\text{pos}, 2i+1)} = \cos\left(\frac{\text{pos}}{10000^{2i/d_{\text{model}}}}\right).
\end{equation}

\subsubsection{Self-Attention Mechanism}

Self-attention allows each token representation to aggregate information from the rest of the sequence. The mechanism uses three learned projections of each input: queries, keys, and values \cite{vaswani2017attention}. Queries represent what a token is looking for, keys represent what each token offers for comparison, and values contain the information that will be aggregated.

For an input sequence representation $\matr{X} \in \reals^{n \times d_{\text{model}}}$, the projections are:
\begin{equation}
    \matr{Q} = \matr{X}\matr{W}^{Q}, \quad
    \matr{K} = \matr{X}\matr{W}^{K}, \quad
    \matr{V} = \matr{X}\matr{W}^{V}.
\end{equation}

Scaled dot-product attention is then computed as:
\begin{equation}
    \operatorname{Attention}(\matr{Q}, \matr{K}, \matr{V})
    =
    \operatorname{softmax}\left(\frac{\matr{Q}\matr{K}^{\top}}{\sqrt{d_k}}\right)\matr{V}.
\end{equation}
The factor $\sqrt{d_k}$ prevents dot products from becoming too large, which would push the softmax into saturated regions with weak gradients.

\subsubsection{Multi-Head Attention}

Multi-head attention applies several attention mechanisms in parallel, allowing the model to attend to different representation subspaces and token relationships at the same time \cite{vaswani2017attention}. For $h$ heads:
\begin{equation}
    \operatorname{head}_i =
    \operatorname{Attention}(\matr{X}\matr{W}^{Q}_{i}, \matr{X}\matr{W}^{K}_{i}, \matr{X}\matr{W}^{V}_{i}),
    \quad i = 1,\ldots,h.
\end{equation}

The head outputs are concatenated and projected back to the model dimension:
\begin{equation}
    \operatorname{MultiHead}(\matr{X}) =
    \operatorname{Concat}(\operatorname{head}_1,\ldots,\operatorname{head}_h)\matr{W}^{O}.
\end{equation}
Here, $\matr{W}^{O} \in \reals^{hd_v \times d_{\text{model}}}$ is the output projection. Attention heads and their projection matrices are important pruning targets because entire heads can sometimes be removed with limited accuracy loss.

\subsubsection{Feed-Forward Network}

After attention mixes information across tokens, a position-wise Feed-Forward Network (FFN) processes each token representation independently:
\begin{equation}
    \operatorname{FFN}(\vect{x}) = \matr{W}_2 \sigma(\matr{W}_1\vect{x} + \vect{b}_1) + \vect{b}_2.
\end{equation}
The FFN usually expands the hidden representation to a larger intermediate dimension $d_{ff}$ and then projects it back to $d_{\text{model}}$. In the original Transformer, $d_{ff}$ is typically much larger than $d_{\text{model}}$, making FFN channels a major source of parameters and a natural target for structured pruning.

\subsubsection{Residual Connections and Layer Normalization}

Each attention or FFN sub-layer is wrapped by a residual path and layer normalization. Residual connections preserve a direct path for information and gradients, while normalization stabilizes the scale of hidden representations. This design is essential for training deep Transformer stacks and is also important when later compression methods remove heads, channels, or blocks \cite{vaswani2017attention,He2016,Ba2016}.

\subsection{Encoder and Decoder Architecture}

With the core components defined, the encoder and decoder can be described cleanly. The full encoder-decoder Transformer contains two main stacks:
\begin{itemize}
    \item \textbf{Encoder:} Processes the complete input sequence in parallel and produces contextual source representations.
    \item \textbf{Decoder:} Generates the output sequence autoregressively while attending to both previous target tokens and the encoder output.
\end{itemize}

The encoder converts the source sequence into a contextual memory representation $\matr{H}_{\text{enc}}$. The decoder uses this memory through cross-attention while also using masked self-attention to prevent access to future target tokens. The complete architecture is shown in Figure~\ref{fig:transformer-architecture}.

\begin{figure}[H]
  \centering
  \begin{tikzpicture}[scale=1.2, transform shape]
    \tikzstyle{Sanode} = [minimum height=1.4em, minimum width=7em, inner sep=3pt, rounded corners=1.5pt, draw=slateblue, fill=mistyblue!30]
    \tikzstyle{Resnode} = [minimum height=1.1em, minimum width=7em, inner sep=3pt, rounded corners=1.5pt, draw=dustyteal, fill=softgray!30]
    \tikzstyle{ffnnode} = [minimum height=1.4em, minimum width=7em, inner sep=3pt, rounded corners=1.5pt, draw=slateblue, fill=sagegreen!30]
    \tikzstyle{outputnode} = [minimum height=1.4em, minimum width=7em, inner sep=3pt, rounded corners=1.5pt, draw=slateblue, fill=mistyblue!50]
    \tikzstyle{inputnode} = [minimum height=1.4em, minimum width=3.5em, inner sep=3pt, rounded corners=1.5pt, draw=dustyteal, fill=mistyblue!20]
    \tikzstyle{posnode} = [minimum height=1.4em, minimum width=3.5em, inner sep=3pt, rounded corners=1.5pt, draw=dustyteal, fill=softgray!50]
    \tikzstyle{standard} = [rounded corners=3pt]

    \node [Sanode, anchor=west] (sa1) at (0,0) {\tiny \textbf{Self-Attention}};
    \node [Resnode, anchor=south] (res1) at ([yshift=0.3em]sa1.north) {\tiny \textbf{Add \& LayerNorm}};
    \node [ffnnode, anchor=south] (ffn1) at ([yshift=1em]res1.north) {\tiny \textbf{FFN}};
    \node [Resnode, anchor=south] (res2) at ([yshift=0.3em]ffn1.north) {\tiny \textbf{Add \& LayerNorm}};
    \node [inputnode, anchor=north west] (input1) at ([yshift=-1em,xshift=-0.5em]sa1.south west) {\tiny \textbf{Emb.}};
    \node [] (add1) at ([yshift=-1.6em,xshift=3.5em]sa1.south west) {$+$};
    \node [posnode, anchor=north east] (pos1) at ([yshift=-1em,xshift=0.5em]sa1.south east) {\tiny \textbf{Position}};
    \node [anchor=south] (encoder) at ([xshift=0.2em,yshift=0.6em]res2.north west) {\scriptsize \textbf{Encoder}};
    \draw [->,dustyteal] ([yshift=-4em]sa1.south) -- (add1);

    \draw [->,dustyteal] (sa1.north) -- (res1.south);
    \draw [->,dustyteal] (res1.north) -- (ffn1.south);
    \draw [->,dustyteal] (ffn1.north) -- (res2.south);
    \draw [->,dustyteal] ([yshift=-1em]sa1.south) -- (sa1.south);

    \node [Sanode, anchor=west] (sa2) at ([xshift=3em]sa1.east) {\tiny \textbf{Masked Self-Attn}};
    \node [Resnode, anchor=south] (res3) at ([yshift=0.3em]sa2.north) {\tiny \textbf{Add \& LayerNorm}};
    \node [Sanode, anchor=south] (ed1) at ([yshift=1em]res3.north) {\tiny \textbf{Cross-Attention}};
    \node [Resnode, anchor=south] (res4) at ([yshift=0.3em]ed1.north) {\tiny \textbf{Add \& LayerNorm}};
    \node [ffnnode, anchor=south] (ffn2) at ([yshift=1em]res4.north) {\tiny \textbf{FFN}};
    \node [Resnode, anchor=south] (res5) at ([yshift=0.3em]ffn2.north) {\tiny \textbf{Add \& LayerNorm}};
    \node [outputnode, anchor=south] (o1) at ([yshift=1em]res5.north) {\tiny \textbf{Linear + Softmax}};
    \node [inputnode, anchor=north west] (input2) at ([yshift=-1em,xshift=-0.5em]sa2.south west) {\tiny \textbf{Emb.}};
    \node [] (add) at ([yshift=-1.6em,xshift=3.5em]sa2.south west) {$+$};
    \node [posnode, anchor=north east] (pos2) at ([yshift=-1em,xshift=0.5em]sa2.south east) {\tiny \textbf{Position}};
    \node [anchor=east] (decoder) at ([xshift=-1em,yshift=-1.5em]o1.west) {\scriptsize \textbf{Decoder}};

    \draw[<-, dustyteal] (add.south) -- ++(0,-1.5em);

    \draw [->,dustyteal] (sa2.north) -- (res3.south);
    \draw [->,dustyteal] (res3.north) -- (ed1.south);
    \draw [->,dustyteal] (ed1.north) -- (res4.south);
    \draw [->,dustyteal] (res4.north) -- (ffn2.south);
    \draw [->,dustyteal] (ffn2.north) -- (res5.south);
    \draw [->,dustyteal] (res5.north) -- (o1.south);
    \draw [->,dustyteal] (o1.north) -- ([yshift=0.5em]o1.north);
    \draw [->,dustyteal] ([yshift=-1em]sa2.south) -- (sa2.south);

    \draw[->,standard,dustyteal] ([yshift=-0.5em]sa1.south) -- ([xshift=-4em,yshift=-0.5em]sa1.south) -- ([xshift=-4em,yshift=2.3em]sa1.south) -- ([xshift=-3.5em,yshift=2.3em]sa1.south);
    \draw[->,standard,dustyteal] ([yshift=0.5em]res1.north) -- ([xshift=-4em,yshift=0.5em]res1.north) -- ([xshift=-4em,yshift=3.3em]res1.north) -- ([xshift=-3.5em,yshift=3.3em]res1.north);
    \draw[->,standard,dustyteal] ([yshift=-0.5em]sa2.south) -- ([xshift=4em,yshift=-0.5em]sa2.south) -- ([xshift=4em,yshift=2.3em]sa2.south) -- ([xshift=3.5em,yshift=2.3em]sa2.south);
    \draw[->,standard,dustyteal] ([yshift=0.5em]res3.north) -- ([xshift=4em,yshift=0.5em]res3.north) -- ([xshift=4em,yshift=3.3em]res3.north) -- ([xshift=3.5em,yshift=3.3em]res3.north);
    \draw[->,standard,dustyteal] ([yshift=0.5em]res4.north) -- ([xshift=4em,yshift=0.5em]res4.north) -- ([xshift=4em,yshift=3.3em]res4.north) -- ([xshift=3.5em,yshift=3.3em]res4.north);
    \draw[->,standard, line width=1.2pt, slateblue] (res2.north) -- ([yshift=0.5em]res2.north) -- ([xshift=5em,yshift=0.5em]res2.north) -- ([xshift=5em,yshift=-2.2em]res2.north) -- ([xshift=6.5em,yshift=-2.2em]res2.north);
    \node[font=\tiny, slateblue] at ([xshift=5.5em,yshift=-1.5em]res2.north) {$\matr{H}_{\text{enc}}$};

  \end{tikzpicture}
  \caption{Complete Encoder-Decoder Transformer architecture. The encoder memory matrix $\matr{H}_{\text{enc}}$ (highlighted in blue) is fed to all decoder layers for cross-attention. Residual connections (curved arrows) enable stable gradient flow throughout the network.}
  \label{fig:transformer-architecture}
\end{figure}

\subsubsection{The Encoder Stack}

An encoder layer contains multi-head self-attention followed by a feed-forward network, with residual connections and normalization around each sub-layer. In encoder self-attention, the queries, keys, and values all come from the same source sequence. This makes the encoder bidirectional: every input token can attend to every other input token. Figure~\ref{fig:encoder} shows the internal order of these operations before the resulting representation is passed to the decoder. After $N$ encoder layers, the output is the memory matrix $\matr{H}_{\text{enc}} \in \reals^{n \times d_{\text{model}}}$, which summarizes the source sequence for the decoder.

\begin{figure}[ht]
  \centering
  \begin{tikzpicture}[scale=1.2, transform shape]
    \tikzstyle{Sanode} = [minimum height=1.4em, minimum width=7em, inner sep=3pt, rounded corners=1.5pt, draw=slateblue, fill=mistyblue!30]
    \tikzstyle{Resnode} = [minimum height=1.1em, minimum width=7em, inner sep=3pt, rounded corners=1.5pt, draw=dustyteal, fill=softgray!30]
    \tikzstyle{ffnnode} = [minimum height=1.4em, minimum width=7em, inner sep=3pt, rounded corners=1.5pt, draw=slateblue, fill=sagegreen!30]
    \tikzstyle{inputnode} = [minimum height=1.4em, minimum width=3.5em, inner sep=3pt, rounded corners=1.5pt, draw=dustyteal, fill=mistyblue!20]
    \tikzstyle{posnode} = [minimum height=1.4em, minimum width=3.5em, inner sep=3pt, rounded corners=1.5pt, draw=dustyteal, fill=softgray!50]
    \tikzstyle{standard} = [rounded corners=3pt]

    \node [Sanode, anchor=west] (sa1) at (0,0) {\tiny \textbf{Self-Attention}};
    \node [Resnode, anchor=south] (res1) at ([yshift=0.3em]sa1.north) {\tiny \textbf{Add \& LayerNorm}};
    \node [ffnnode, anchor=south] (ffn1) at ([yshift=1em]res1.north) {\tiny \textbf{FFN}};
    \node [Resnode, anchor=south] (res2) at ([yshift=0.3em]ffn1.north) {\tiny \textbf{Add \& LayerNorm}};
    \node [inputnode, anchor=north west] (input1) at ([yshift=-1em,xshift=-0.5em]sa1.south west) {\tiny \textbf{Emb.}};
    \node [] (add1) at ([yshift=-1.6em,xshift=3.5em]sa1.south west) {$+$};
    \node [posnode, anchor=north east] (pos1) at ([yshift=-1em,xshift=0.5em]sa1.south east) {\tiny \textbf{Position}};
    \node [anchor=south] (encoder) at ([xshift=0.2em,yshift=0.6em]res2.north west) {\scriptsize \textbf{Encoder}};

    \draw[<-, dustyteal] (add1.south) -- ++(0,-1.5em);

    \draw [->,dustyteal] (sa1.north) -- (res1.south);
    \draw [->,dustyteal] (res1.north) -- (ffn1.south);
    \draw [->,dustyteal] (ffn1.north) -- (res2.south);
    \draw [->,dustyteal] ([yshift=-1em]sa1.south) -- (sa1.south);

    \draw[->,standard,dustyteal] ([yshift=-0.5em]sa1.south) -- ([xshift=-4em,yshift=-0.5em]sa1.south) -- ([xshift=-4em,yshift=2.3em]sa1.south) -- ([xshift=-3.5em,yshift=2.3em]sa1.south);
    \draw[->,standard,dustyteal] ([yshift=0.5em]res1.north) -- ([xshift=-4em,yshift=0.5em]res1.north) -- ([xshift=-4em,yshift=3.3em]res1.north) -- ([xshift=-3.5em,yshift=3.3em]res1.north);

  \end{tikzpicture}
  \caption{Transformer encoder layer architecture. The input embedding (Emb.) is combined with positional information, processed by self-attention and a Feed-Forward Network (FFN), and stabilized by residual connections with Layer Normalization (LayerNorm).}
  \label{fig:encoder}
\end{figure}

\subsubsection{The Decoder Stack}

A decoder layer adds masked self-attention and cross-attention before the feed-forward network, enabling autoregressive generation while attending to encoder outputs. Masked self-attention restricts token $t$ to positions $\leq t$, preventing information leakage from future target tokens. Cross-attention then links the decoder to the encoder: decoder states act as queries, while encoder states provide keys and values. Figure~\ref{fig:decoder} shows this additional cross-attention block, which is the main structural difference between the decoder layer and the encoder layer.

\begin{figure}[htbp]
  \centering
  \begin{tikzpicture}[scale=1.2, transform shape]
    \tikzstyle{Sanode} = [minimum height=1.4em, minimum width=7em, inner sep=3pt, rounded corners=1.5pt, draw=slateblue, fill=mistyblue!30]
    \tikzstyle{Resnode} = [minimum height=1.1em, minimum width=7em, inner sep=3pt, rounded corners=1.5pt, draw=dustyteal, fill=softgray!30]
    \tikzstyle{ffnnode} = [minimum height=1.4em, minimum width=7em, inner sep=3pt, rounded corners=1.5pt, draw=slateblue, fill=sagegreen!30]
    \tikzstyle{outputnode} = [minimum height=1.4em, minimum width=7em, inner sep=3pt, rounded corners=1.5pt, draw=slateblue, fill=mistyblue!50]
    \tikzstyle{inputnode} = [minimum height=1.4em, minimum width=3.5em, inner sep=3pt, rounded corners=1.5pt, draw=dustyteal, fill=mistyblue!20]
    \tikzstyle{posnode} = [minimum height=1.4em, minimum width=3.5em, inner sep=3pt, rounded corners=1.5pt, draw=dustyteal, fill=softgray!50]
    \tikzstyle{standard} = [rounded corners=3pt]

    \node [Sanode, anchor=west] (sa2) at (0,0) {\tiny \textbf{Masked Self-Attn}};
    \node [Resnode, anchor=south] (res3) at ([yshift=0.3em]sa2.north) {\tiny \textbf{Add \& LayerNorm}};
    \node [Sanode, anchor=south] (ed1) at ([yshift=1em]res3.north) {\tiny \textbf{Cross-Attention}};
    \node [Resnode, anchor=south] (res4) at ([yshift=0.3em]ed1.north) {\tiny \textbf{Add \& LayerNorm}};
    \node [ffnnode, anchor=south] (ffn2) at ([yshift=1em]res4.north) {\tiny \textbf{FFN}};
    \node [Resnode, anchor=south] (res5) at ([yshift=0.3em]ffn2.north) {\tiny \textbf{Add \& LayerNorm}};
    \node [outputnode, anchor=south] (o1) at ([yshift=1em]res5.north) {\tiny \textbf{Linear + Softmax}};
    \node [inputnode, anchor=north west] (input2) at ([yshift=-1em,xshift=-0.5em]sa2.south west) {\tiny \textbf{Emb.}};
    \node [] (add) at ([yshift=-1.6em,xshift=3.5em]sa2.south west) {$+$};
    \node [posnode, anchor=north east] (pos2) at ([yshift=-1em,xshift=0.5em]sa2.south east) {\tiny \textbf{Position}};
    \node [anchor=south] (decoder) at ([xshift=0.2em,yshift=0.6em]o1.north west) {\scriptsize \textbf{Decoder}};

    \draw [->,dustyteal] ([yshift=-4em]sa2.south) -- (add);

    \draw [->,dustyteal] (sa2.north) -- (res3.south);
    \draw [->,dustyteal] (res3.north) -- (ed1.south);
    \draw [->,dustyteal] (ed1.north) -- (res4.south);
    \draw [->,dustyteal] (res4.north) -- (ffn2.south);
    \draw [->,dustyteal] (ffn2.north) -- (res5.south);
    \draw [->,dustyteal] (res5.north) -- (o1.south);
    \draw [->,dustyteal] (o1.north) -- ([yshift=0.5em]o1.north);
    \draw [->,dustyteal] ([yshift=-1em]sa2.south) -- (sa2.south);

    \draw[->,standard,dustyteal] ([yshift=-0.5em]sa2.south) -- ([xshift=4em,yshift=-0.5em]sa2.south) -- ([xshift=4em,yshift=2.3em]sa2.south) -- ([xshift=3.5em,yshift=2.3em]sa2.south);
    \draw[->,standard,dustyteal] ([yshift=0.5em]res3.north) -- ([xshift=4em,yshift=0.5em]res3.north) -- ([xshift=4em,yshift=3.3em]res3.north) -- ([xshift=3.5em,yshift=3.3em]res3.north);
    \draw[->,standard,dustyteal] ([yshift=0.5em]res4.north) -- ([xshift=4em,yshift=0.5em]res4.north) -- ([xshift=4em,yshift=3.3em]res4.north) -- ([xshift=3.5em,yshift=3.3em]res4.north);

  \end{tikzpicture}
  \caption{Transformer decoder layer architecture. The decoder contains masked self-attention for previous target tokens, cross-attention for reading encoder outputs, and a Feed-Forward Network (FFN), each wrapped with residual connections and Layer Normalization (LayerNorm).}
  \label{fig:decoder}
\end{figure}

\subsubsection{Encoder-Decoder Information Flow}

The complete sequence-to-sequence computation proceeds as follows:
\begin{enumerate}[leftmargin=*]
  \item \textbf{Encoding phase.} Source tokens $\vect{x}=(x_1,\ldots,x_n)$ are embedded, combined with positional information, and passed through $N$ encoder layers. The output is $\matr{H}_{\text{enc}}$.
  \item \textbf{Decoding phase.} Target tokens $\vect{y}=(y_1,\ldots,y_t)$ are embedded and processed autoregressively. Each decoder layer applies masked self-attention, cross-attention over $\matr{H}_{\text{enc}}$, and a feed-forward transformation.
  \item \textbf{Output phase.} The final decoder representation is projected to vocabulary logits. A softmax layer converts logits into token probabilities, and the selected token is fed back into the decoder at the next step.
\end{enumerate}

\vspace{0.8em}
\noindent\rule{\textwidth}{0.4pt}
\vspace{0.8em}

At this point in the report, we can view the Transformer as an arrangement of embedding layers, attention layers, feed-forward functions, and encoder--decoder connections. This architecture accounts for why it performs well on sequence modelling tasks, since each position can refine its representation using information from the rest of the sequence. This expressive structure, however, is computationally demanding. The next point examines where this cost arises inside the Transformer architecture.

\subsection{The Computational Bottleneck of Transformers}
\label{sec:transformer_bottleneck}
While powerful, the Transformer's success comes at a significant computational cost, which forms a central motivation for the compression techniques explored in this report. The standard Transformer architecture replaces recurrence with self-attention, enabling strong parallelism and high performance across sequence modelling tasks~\cite{vaswani2017attention}. However, this advantage comes with a major bottleneck: in full self-attention, each token attends to every other token in the sequence. For a sequence of length $n$, this requires an attention score matrix of size $n \times n$, making the attention component scale quadratically with sequence length in both computation and memory~\cite{beltagy2020longformer,tay2020efficient}.

This quadratic behaviour has direct practical implications. Doubling the sequence length from $512$ to $1024$ multiplies the attention matrix size by four, since $1024^2 / 512^2 = 4$. As a result, processing long documents or long-context inputs becomes substantially more expensive, which has motivated a large body of work on efficient Transformer variants and sparse or linear attention mechanisms~\cite{beltagy2020longformer,tay2020efficient}.

At the same time, modern language models benefit strongly from scaling model size, training data, and compute~\cite{kaplan2020scaling}. Later scaling-law work further showed that model size and the number of training tokens should be scaled together under a fixed compute budget, highlighting the continuing importance of large-scale training~\cite{hoffmann2022training}. This creates a tension between scale and deployment: larger models often achieve better performance, but practical inference is constrained by latency, memory, energy, and hardware availability. This tension motivates compression methods such as pruning, quantization, and knowledge distillation~\cite{han2015deep,sanh2019distilbert}.

\par\smallskip
\noindent\rule{0.35\linewidth}{0.4pt}
\par\smallskip

The computational bottleneck defined above can be considered in several different ways. A Transformer may be expensive because of the size of its parameters, the memory required for activations, the number of arithmetic operations it performs, or the time required to run it on a specific hardware platform. Before discussing model efficiency in more detail, it is therefore useful to introduce the most common measures used to assess computational cost.

\subsection{Hardware-Aware Efficiency Metrics}

Compression is evaluated through several distinct efficiency metrics. Parameter count, memory footprint, arithmetic cost, latency, and throughput capture different aspects of model execution. Hardware-aware evaluation therefore requires separate treatment of storage cost, arithmetic cost, and runtime behavior~\cite{Sze2017EfficientDNN}.

\subsubsection{Parameter Count and Model Size}

Let $\params = \{W^{(l)}, b^{(l)}\}_{l=1}^{L}$ denote the parameters of a network. The total parameter count is
\begin{equation}
P = \sum_{l=1}^{L} \left( |W^{(l)}| + |b^{(l)}| \right).
\end{equation}
If the parameters are stored with precision $b_w$ bits, the static model footprint is
\begin{equation}
M_{\text{model}} = \frac{P\, b_w}{8} \quad \text{bytes}.
\end{equation}
This follows from counting the stored elements and multiplying by the number of bytes used by each element~\cite{PyTorchNumel,PyTorchElementSize}.
This quantity determines whether a model fits within device memory and how much parameter traffic must be sustained during inference. Quantization reduces $b_w$, pruning reduces the number of active parameters, and low-rank factorization reduces the parameterization itself.

For sparse models, storage depends on both the remaining weights and the associated metadata. If $P_{\text{nz}}$ denotes the number of non-zero weights and $M_{\text{meta}}$ the storage needed for indices or masks, then
\begin{equation}
M_{\text{sparse}} = \frac{P_{\text{nz}}\, b_w}{8} + M_{\text{meta}}.
\end{equation}

\subsubsection{MACs, FLOPs, and Arithmetic Intensity}

For a dense matrix multiplication $Y = AB$ with $A \in \reals^{m \times k}$ and $B \in \reals^{k \times n}$, the arithmetic cost is
\begin{equation}
\mathrm{MACs}(A,B) = mkn,
\qquad
\mathrm{FLOPs}(A,B) \approx 2mkn.
\end{equation}
The distinction is conventional: one multiply-accumulate corresponds to one MAC, while FLOP counting often treats the multiply and the addition as two floating-point operations~\cite{NVIDIAMatrixMultiplication}. In Transformer layers, the feed-forward subnetwork contributes large dense matrix products, while self-attention adds the sequence-dependent $O(n^2 d_{\text{model}})$ cost of score computation and value aggregation.

Arithmetic cost alone is insufficient for hardware analysis. A useful complementary quantity is the arithmetic intensity
\begin{equation}
I = \frac{\mathrm{FLOPs}}{\mathrm{Bytes\ Moved}},
\end{equation}
which measures how much computation is performed per byte transferred~\cite{Williams2009Roofline,NVIDIAMatrixMultiplication}. Low arithmetic intensity indicates a memory-bound workload, whereas high arithmetic intensity indicates a compute-bound workload.

\subsubsection{Bandwidth, Latency, and Throughput}

Memory bandwidth measures the rate at which data can be transferred between memory and compute units. If $B_{\text{move}}$ bytes are transferred during an operation that takes time $T$, the effective bandwidth demand is
\begin{equation}
\mathrm{BW}_{\text{eff}} = \frac{B_{\text{move}}}{T}.
\end{equation}
This follows the standard effective-bandwidth definition used for GPU kernels~\cite{CUDABestPractices}.
When parameter streaming and activation movement dominate execution, bandwidth becomes the primary runtime constraint.

Latency and throughput characterize runtime from complementary viewpoints. Latency is the time required to complete one inference request or one decoding step~\cite{Sze2017EfficientDNN}:
\begin{equation}
T_{\text{latency}} = T_{\text{compute}} + T_{\text{memory}} + T_{\text{overhead}}.
\end{equation}
Throughput measures completed work per unit time, for example requests per second or generated tokens per second~\cite{Sze2017EfficientDNN}:
\begin{equation}
\mathrm{Throughput} = \frac{N_{\text{outputs}}}{T}.
\end{equation}
In large autoregressive models, latency and throughput often follow different scaling trends. Batching can increase throughput while single-request latency remains high. Parameter compression can reduce memory transfer while leaving runtime gains limited when execution is poorly aligned with the hardware schedule.

% !TEX root = ../../main_without_quantization.tex
\section{Genetic Algorithms and Multi-Objective Search}
\label{sec:multiobjective_genetic_algorithms}

With the hardware-aware metrics introduced above, it becomes possible to compare model variants concretely. In particular, different designs may be compared with regard to prediction accuracy, parameter count, memory footprint, computational cost, memory bandwidth usage, and latency/throughput. In compression, these variants may correspond to different pruning masks, retained heads, feed-forward widths, removed layers, or layer-wise numerical precision.

Taken together, these choices form a large design space. Since enumerating all possible candidates is usually prohibitive, a search procedure is needed to explore this space and identify candidates that satisfy efficiency requirements while preserving prediction accuracy.

To manage this search, one can maintain a set of candidate designs instead of tracking a single design at a time. Each candidate corresponds to a possible compressed model. The candidates are evaluated, the more relevant ones are kept with higher probability, and new candidates are generated from existing ones through small changes or recombination. By repeating this process, the candidate set can move toward desirable regions of the design space.

A population-based search algorithm of this kind is known as a genetic algorithm. In this terminology, a candidate compressed model is an individual, the set of candidates is a population, and the evaluation measure is the fitness function. The following section introduces these notions together with encoding, selection, crossover, and mutation, then links them to compression design choices, multi-objective optimization, and NSGA-II \cite{holland1975adaptation,goldberg1989genetic,deb2001multi}.

\subsection{Genetic Algorithm Fundamentals}

A Genetic Algorithm (GA) is a population-based search method inspired by biological evolution. It maintains candidate solutions, represents them in a form that can be modified, evaluates their quality, and repeatedly generates new candidates through selection, crossover, and mutation \cite{holland1975adaptation,goldberg1989genetic,michalewicz1996genetic}.

\subsubsection{Components of a Genetic Algorithm}

The basic object manipulated by a GA is an individual, also called a chromosome. Each individual encodes one candidate solution to the problem. A population is then a set of such individuals evaluated during the same generation. The fitness function assigns a quality score to each candidate, and the genetic operators define how new candidates are produced from existing ones \cite{goldberg1989genetic,michalewicz1996genetic}. These definitions establish what the algorithm manipulates; the next issue is how a candidate is written as a chromosome.

In model compression, an individual may represent a compressed architecture. For example, consider the chromosome
\begin{equation}
	\mathbf{c} = [1,0,1,1 \mid 2048,1024 \mid 8,4,4,8],
\end{equation}
where the first block indicates which attention heads are retained, the second block gives selected feed-forward widths, and the final block assigns layer-wise bit-widths. The fitness value then measures how good this candidate is with respect to the desired objective, such as validation accuracy, parameter count, memory footprint, or latency.

\subsubsection{Encoding Methods for Individuals}

The encoding determines how a candidate solution is written as a chromosome and therefore what kind of changes the algorithm can apply \cite{michalewicz1996genetic}. A \textbf{binary encoding} represents each gene with one bit, so every decision has two possible states. In compression, a bit value of $1$ may indicate that a neuron, head, or channel is retained, while $0$ indicates that it is removed. An \textbf{integer encoding} represents each gene with an integer selected from a finite set of choices. This is useful when a layer must choose one option among several possible bit-widths, such as $\{2,4,8,16\}$.

The choice of encoding is important because it determines what the genetic operators can change. A poor encoding may create invalid candidates or make useful solutions hard to reach. A good encoding reflects the structure of the compression problem while keeping the search space manageable. Once this representation is fixed, the algorithm can operate on a population across generations.

\subsubsection{Operation of Genetic Algorithms}

A GA begins by initializing a population, either randomly or from existing candidate designs. Each individual is evaluated using the fitness function. The algorithm then selects promising candidates, recombines their encoded decisions through crossover, and applies mutation to preserve exploration \cite{goldberg1989genetic}. The resulting candidates form the next population, and this cycle is repeated until a stopping criterion is reached, such as a maximum number of generations, convergence of the population, or an evaluation budget. The three operators below define how this update takes place.

\paragraph{Selection Operator.}
Selection chooses which individuals are allowed to reproduce. Better candidates should have a higher chance of being selected, but weaker candidates are not always discarded immediately because they may contain useful partial structures. Common strategies include roulette-wheel selection, rank-based selection, and tournament selection \cite{baker1985adaptive,miller1995genetic}. The selected candidates then serve as parents for crossover.

\paragraph{Crossover Operator.}
Crossover combines information from two or more parent individuals to produce offspring. In binary or integer encodings, this may exchange segments of chromosomes between parents. In uniform crossover, each gene can be inherited independently from either parent \cite{syswerda1989uniform}. For compression, crossover can combine useful pruning patterns or precision assignments discovered in different candidates. Since crossover reuses information already present in the population, mutation is used to introduce local changes.

\paragraph{Mutation Operator.}
Mutation introduces small random changes into an individual. It prevents the population from collapsing too early around one region of the search space and allows new structures to appear \cite{back1993optimal}. In a compression chromosome, mutation may flip a keep-or-remove bit, change the bit-width of a layer, or adjust a pruning threshold.

\vspace{0.4em}
\noindent\rule{0.35\textwidth}{0.4pt}
\vspace{0.4em}

With these mechanics defined, the compression problem can now be written as a genetic search problem.

\subsection{Compression as a Search Problem}

With the GA terminology in place, compression can be described as a search over encoded model variants. Let $\bx$ denote a compression configuration, such as a pruning mask, a vector of layer-wise bit-widths, or a set of structural choices. Evaluating $\bx$ produces a compressed model whose quality depends on both predictive performance and deployment cost.

A single scalar score is sometimes used, for example by combining accuracy loss and model size into one weighted objective. However, this requires choosing the trade-off weights before the search begins. In practice, compression often requires a set of alternatives: one model may preserve more accuracy, another may use less memory, and another may reduce latency more aggressively. This need for several valid trade-offs leads to multi-objective optimization and Pareto dominance.

\subsection{Multi-Objective Optimization and Pareto Dominance}

\subsubsection{Problem Formulation}
Multi-objective optimization formalizes the comparison introduced above by replacing one scalar objective with an objective vector. A Multi-Objective Optimization Problem (MOP) minimizes a vector of objective functions \cite{deb2001multi}:

\begin{equation}
	\begin{aligned}
		& \underset{\bx}{\text{minimize}}
		& & \bfun(\bx) = [f_1(\bx), f_2(\bx), \dots, f_M(\bx)]^T \\
		& \text{subject to}
		& & g_j(\bx) \leq 0, \quad j = 1, \dots, J \\
		& & & x_i^{(L)} \leq x_i \leq x_i^{(U)}, \quad i = 1, \dots, n .
	\end{aligned}
\end{equation}

Here, $\bx \in \mathbb{R}^n$ is the decision vector and $\bfun: \mathbb{R}^n \to \mathbb{R}^M$ maps each candidate to the objective space. In compression, the objectives may include validation error, number of parameters, model footprint, FLOPs, latency, or energy consumption.

Since each candidate is represented by several objective values, candidates cannot always be ranked by a single numerical score. The comparison is therefore based on whether one candidate improves at least one objective without worsening the others. This idea is formalized by Pareto dominance.

\subsubsection{Pareto Dominance}

\paragraph{Definition (Pareto Dominance).}
A vector $\mathbf{u}$ dominates $\mathbf{v}$, denoted $\mathbf{u} \prec \mathbf{v}$, if and only if:
\begin{equation}
	\forall i: u_i \leq v_i \quad \land \quad \exists j: u_j < v_j .
\end{equation}

Here, $\mathbf{u}$ and $\mathbf{v}$ are objective vectors, and $u_i$ and $v_i$ are their values on objective $i$. Thus, $\mathbf{u}$ is no worse than $\mathbf{v}$ in every objective and strictly better in at least one objective.

\paragraph{Definition (Pareto Optimal Set).}
The Pareto optimal set in decision space is:
\begin{equation}
	PS^* = \{ \bx \in \Omega \mid \nexists \by \in \Omega: \bfun(\by) \prec \bfun(\bx) \}.
\end{equation}

Here, $\Omega$ is the feasible search space, $\bx$ and $\by$ are candidate solutions, and $\bfun(\by) \prec \bfun(\bx)$ means that candidate $\by$ dominates candidate $\bx$ in objective space. Therefore, $PS^*$ contains the feasible candidates that are not dominated by any other feasible candidate.

The corresponding Pareto front in objective space is:
\begin{equation}
	PF^* = \{ \bfun(\bx) \mid \bx \in PS^* \}.
\end{equation}

Here, $PF^*$ is the image of the Pareto optimal set under the objective function $\bfun$. For model compression, the Pareto front gives a set of non-dominated compressed models. A user can then choose the model that best matches the deployment constraint, rather than relying on a single fixed trade-off. Pareto dominance provides the ranking principle, and NSGA-II combines this ranking with the population-based mechanics of genetic algorithms.

\subsection{NSGA-II}

NSGA-II, the Non-dominated Sorting Genetic Algorithm II, is a widely used multi-objective genetic algorithm that connects GA search with Pareto-based ranking \cite{deb2002fast}. It keeps the population-based structure of a GA, but replaces scalar fitness ranking with Pareto-based selection. This makes it suitable for compression problems where accuracy, memory, and runtime objectives must be considered together.

The first key mechanism is non-dominated sorting. Candidates are grouped into fronts: the first front contains candidates that are not dominated by any other candidate, the second front contains candidates dominated only by members of the first front, and so on. This gives priority to candidates closer to the Pareto front.

The second mechanism is crowding distance. Within the same front, NSGA-II favors candidates located in less crowded regions of the objective space. This preserves diversity and prevents the search from returning many nearly identical compression configurations.

The third mechanism is elitist preservation. Parents and offspring are merged, ranked by non-dominated sorting, and filtered using crowding distance when a front cannot fit entirely into the next generation. As a result, strong candidates are preserved while the population continues exploring new trade-offs.

In the context of this report, NSGA-II provides the bridge between genetic search and hardware-aware compression. The chromosome defines the compression decision, the objectives measure predictive and deployment cost, and the algorithm searches for a Pareto set of compressed Transformer variants.

This completes the foundations needed for the proposed contribution: the model can be measured through hardware-aware metrics, encoded as a candidate design, and searched under several objectives at once.

%


\section{Conclusion}

This chapter established the mathematical and architectural foundations required for the study of compression methods in the following chapters. It introduced the neuron model, activation functions, multilayer perceptrons, learning objectives, gradient-based optimization, and backpropagation, then surveyed representative neural architectures, including recurrent, autoencoding, and Transformer-based models.

\noindent The discussion of hardware-aware efficiency formalized the principal deployment metrics that later compression methods seek to improve, including parameter count, memory footprint, arithmetic cost, bandwidth demand, latency, and throughput. The chapter then introduced multi-objective genetic algorithms as a search framework for problems involving competing objectives and discrete design variables. Together, these foundations motivate the compression formulations developed later in the report, where predictive quality must be balanced against memory, runtime, and hardware constraints.

\noindent Within this conceptual frame, pruning emerges as a natural next step. Once neural computation is expressed through parameterized operators, decomposed into identifiable architectural units, and evaluated through resource-sensitive metrics, the central pruning question becomes precise: which parts of the network may be removed with limited functional loss, and under what deployment constraints does that removal become meaningful.
